% Options for packages loaded elsewhere
\PassOptionsToPackage{unicode}{hyperref}
\PassOptionsToPackage{hyphens}{url}
\documentclass[
]{article}
\usepackage{xcolor}
\usepackage[margin=1in]{geometry}
\usepackage{amsmath,amssymb}
\setcounter{secnumdepth}{-\maxdimen} % remove section numbering
\usepackage{iftex}
\ifPDFTeX
  \usepackage[T1]{fontenc}
  \usepackage[utf8]{inputenc}
  \usepackage{textcomp} % provide euro and other symbols
\else % if luatex or xetex
  \usepackage{unicode-math} % this also loads fontspec
  \defaultfontfeatures{Scale=MatchLowercase}
  \defaultfontfeatures[\rmfamily]{Ligatures=TeX,Scale=1}
\fi
\usepackage{lmodern}
\ifPDFTeX\else
  % xetex/luatex font selection
\fi
% Use upquote if available, for straight quotes in verbatim environments
\IfFileExists{upquote.sty}{\usepackage{upquote}}{}
\IfFileExists{microtype.sty}{% use microtype if available
  \usepackage[]{microtype}
  \UseMicrotypeSet[protrusion]{basicmath} % disable protrusion for tt fonts
}{}
\makeatletter
\@ifundefined{KOMAClassName}{% if non-KOMA class
  \IfFileExists{parskip.sty}{%
    \usepackage{parskip}
  }{% else
    \setlength{\parindent}{0pt}
    \setlength{\parskip}{6pt plus 2pt minus 1pt}}
}{% if KOMA class
  \KOMAoptions{parskip=half}}
\makeatother
\usepackage{graphicx}
\makeatletter
\newsavebox\pandoc@box
\newcommand*\pandocbounded[1]{% scales image to fit in text height/width
  \sbox\pandoc@box{#1}%
  \Gscale@div\@tempa{\textheight}{\dimexpr\ht\pandoc@box+\dp\pandoc@box\relax}%
  \Gscale@div\@tempb{\linewidth}{\wd\pandoc@box}%
  \ifdim\@tempb\p@<\@tempa\p@\let\@tempa\@tempb\fi% select the smaller of both
  \ifdim\@tempa\p@<\p@\scalebox{\@tempa}{\usebox\pandoc@box}%
  \else\usebox{\pandoc@box}%
  \fi%
}
% Set default figure placement to htbp
\def\fps@figure{htbp}
\makeatother
\setlength{\emergencystretch}{3em} % prevent overfull lines
\providecommand{\tightlist}{%
  \setlength{\itemsep}{0pt}\setlength{\parskip}{0pt}}
\usepackage{tikz}
\usepackage{pgfplots}
\usepackage{xcolor}
\usepackage{bookmark}
\IfFileExists{xurl.sty}{\usepackage{xurl}}{} % add URL line breaks if available
\urlstyle{same}
\hypersetup{
  hidelinks,
  pdfcreator={LaTeX via pandoc}}

\author{}
\date{\vspace{-2.5em}}

\begin{document}

\section{Question}\label{question}

La grafica muestra informacion de las poblaciones de 5 paises desde 1980
hasta 2023.

\includegraphics[width=14cm,height=\textheight,keepaspectratio]{grafica_poblaciones_paises_pgfplots.png}

\hfill\break
Los datos representan la evolucion demografica de diferentes naciones a
lo largo de varias decadas, mostrando distintas tendencias de
crecimiento poblacional.

Aproximadamente, ¿en que ano las poblaciones de \textbf{Canada} y
\textbf{Finlandia} fueron iguales?

\subsection{Answerlist}\label{answerlist}

\begin{itemize}
\tightlist
\item
  1985
\item
  1990
\item
  1995
\item
  2023
\end{itemize}

\section{Solution}\label{solution}

Para resolver este problema, debemos identificar el punto donde las
lineas de \textbf{Canada} y \textbf{Finlandia} se intersectan en la
grafica.

\textbf{Analisis visual de las tendencias:}

\begin{itemize}
\tightlist
\item
  \textbf{Canada}: Inicia con una poblacion menor pero presenta un
  crecimiento mas acelerado.
\item
  \textbf{Finlandia}: Comienza con una poblacion mayor pero su
  crecimiento es mas moderado.
\end{itemize}

\textbf{Identificacion del punto de interseccion:}

Observando cuidadosamente la grafica, las dos lineas se cruzan
aproximadamente en el ano \textbf{1995}, donde ambas poblaciones
alcanzan un valor similar de aproximadamente 320.000.000 habitantes.

\textbf{Verificacion:}

\begin{itemize}
\tightlist
\item
  Antes de 1995: Finlandia \textgreater{} Canada
\item
  Despues de 1995: Canada \textgreater{} Finlandia
\end{itemize}

\subsection{Answerlist}\label{answerlist-1}

\begin{itemize}
\tightlist
\item
  \textbf{Falso}. El ano 1985 no corresponde al punto de interseccion de
  las lineas.
\item
  \textbf{Falso}. El ano 1990 no corresponde al punto de interseccion de
  las lineas.
\item
  \textbf{Verdadero}. El ano 1995 corresponde al punto de interseccion
  visual entre las lineas de Canada y Finlandia .
\item
  \textbf{Falso}. El ano 2023 no corresponde al punto de interseccion de
  las lineas.
\end{itemize}

\section{Meta-information}\label{meta-information}

exname: Interseccion Poblaciones Paises Grafica Lineas PGFPlots extype:
schoice exsolution: 0010 exshuffle: TRUE exsection:
Estadistica/Interpretacion de Graficas

\end{document}
